\section{AFS 제안과 피드백}
\begin{table}[H]
\centering\small
\caption{AFS에 제공하는 행동 관측}
\begin{tabularx}{\linewidth}{@{}p{38mm}Y@{}}
\toprule
관측 정보 & 활용 \\
\midrule
목표점 경로 거절·반복 & 같은 지점에 대한 거절 반복 파악 \\
중간 목표 도착 & 부분 이동 성공과 전체 과제 완료의 차이 파악 \\
목표 거리 변화·최종 상태 & 실제 진전과 종료 상태 구분 \\
조작 요청과 그 결과 & 요청과 실제 물리 실행의 구분 \\
접촉 시간·물체 이동량 (로봇팔) & 접근 차단과 운반·놓기 실패의 구분 \\
예상 유형과 측정 유형 & 예상이 실제 검출로 이어졌는지 확인 \\
\bottomrule
\end{tabularx}
\end{table}

\textbf{제안 예시 (로봇팔, 시드 661, 탐색 1회차).} 가설: ``그릇에서 접시로 옮기는 경로 중앙에
더 크고 높은 상자를 두면 로봇의 움직임이 막혀 그릇을 옮기지 못할 것이다.''
제안값: 그릇 (-0.080, 0.007), 접시 (0.084, 0.239), 장애물 위치 (-0.041, 0.118), 폭·깊이 0.06\,m,
높이 0.17\,m, 회전 0. 결과: FAIL, 로봇–장애물 접촉 11.0초.
